Este capítulo establece las bases operativas y conceptuales de la investigación aplicadas al sector financiero de la economía social y solidaria. A partir de la experiencia práctica, se identifican limitantes críticos en la gestión de datos crediticios mediante tecnologías modernas como el aprendizaje automático. Finalmente, se establece la hipótesis, los objetivos y las contribuciones que espero lograr con la solución que propongo.

\section{"Background" y Contexto}

Las instituciones financieras, en especial las pertenecientes al sector de la economía social y solidaria, manejan grandes volúmenes de datos crediticios que, correctamente analizados, pueden revelar patrones predictivos sobre el comportamiento de pago de sus socios y clientes. \\
Tradicionalmente, el análisis de cartera se realiza mediante indicadores retrospectivos —tasa de morosidad, índice de castigos, relación cartera vencida / cartera total— que informan sobre lo que ya ocurrió, pero no ofrecen capacidad de anticipación.\\
En los últimos años, la orquestación de las tres tecnologías ha abierto una nueva frontera  de desarrollo:
\\
\newline\textbf{Las Arquitecturas de datos modernas} como los Data Warehouse's y los Data Marts permiten consolidar datos operativos y predicciones en un solo repositorio optimizado para consultas de negocio, eliminando la fricción entre los equipos de datos y los usuarios de negocio.
\\
\newline\textbf{El Aprendizaje profundo} y los modelos de Redes Neuronales Convolucionales han mostrado capacidad para extraer patrones temporales en series financieras, superando a métodos estadísticos clásicos como ARIMA o la regresión logística en tareas de clasificación de riesgo.

\section{La Motivación}

\begin{flushright}
    \small\itshape 
    "En temas de datos, la falta de acceso es tan útil como no tenerlos" 
    \vspace{0.2cm}  
    \normalfont--- Autoría propia. 
\end{flushright}

La motivación de este trabajo surge de la experiencia práctica en el área de desarrollo y el análisis de datos en una institución financiera ecuatoriana; y, al cursar la maestría en inteligencia artificial nos ha permitido identificar las siguientes problemáticas recurrentes:

\begin{itemize}
    \item Los reportes de cartera se generan de forma manual, consumiendo días de trabajo cada mes.
    \item  Los modelos predictivos existentes tienen alta exactitud (accuracy > 93\% ) pero precisión y recall en cero, lo que indica que no detectan crisis reales (problema de clases desbalanceadas).
    \item  No existe un datamart unificado que permita a los analistas cruzar datos históricos con predicciones del modelo.
    \item  Los dashboards son estáticos, dependen de personal técnico para su actualización y su acceso está restringido a un grupo reducido de analistas, lo que contradice la política de puertas abiertas de la institución y limita la democratización de la información.
    \item  No se tiene una herramienta predictiva para mejorar la colocación de capitales en prestamos, lo que podría reducir la tasa de morosidad y mejorar los rendimientos de la cartera.   
\end{itemize}

\section{Planteamiento del problema}

Luego de más de una década desempeñándome en instituciones financieras de la economía social y solidaria, he constatado que, si bien a diario se generan millones de registros transaccionales, su valor estratégico se diluye debido a cuatro brechas estructurales:

\begin{enumerate}
    \item \textbf{Falta de integración de la IA con los sistemas de reporte}: los modelos y técnicas de inteligencia artificial no están disponibles para la construcción de datamarts, lo que impide que los analistas de negocio aprovechen su potencial analítico.
    \item \textbf{Los dashboards no son dinámicos ni auto-gestionados}: La creación y actualización de dashboards requiere intervención manual de personal técnico, creando cuellos de botella.
    \item \textbf{No existe un canal conversacional para consultar los datos}: Los usuarios de negocio deben depender del equipo de datos para responder preguntas que podrían resolverse con lenguaje natural, requiere de interpretación técnica.
    \item \textbf{El monitoreo del modelo no está integrado con las alertas de negocio}: No hay un mecanismo que notifique automáticamente cuando las predicciones superan umbrales críticos o que indiquen un riesgo crediticio.
\end{enumerate}

Estas limitaciones derivan en una  \textbf{Pregunta de investigación}: ¿Cómo integrar un modelo predictivo de riesgo crediticio con un datamart analítico y dashboards, a fin de conformar una plataforma unificada de inteligencia de negocio que permita controlar eficazmente la mora?

\section{Hipótesis}

Los modelos de Gradient Boosting superan a las redes neuronales convolucionales y perceptrones multicapa en la predicción de incumplimiento crediticio multi-horizonte cuando se trabaja con datos tabulares de instituciones de economía social y solidaria, caracterizados por secuencias temporales cortas, clases desbalanceadas y volumen de datos limitado.

\section{Objetivo General}

Desarrollar una plataforma de inteligencia de negocio crediticia que, mediante un flujo de operaciones de Machine Learning automatizado, compare modelos de deep learning y gradient boosting para la predicción multi-horizonte de riesgo, e integre el modelo seleccionado en un datamart analítico con dashboards interactivos para el control de la mora.

\section{Objetivos Específicos}

\begin{enumerate}
    \item  \textbf{Identificar las causas del fracaso de las redes neuronales} en la predicción de incumplimiento crediticio con datos tabulares de instituciones de economía social y solidaria.
    \item \textbf{Determinar qué características de los datos crediticios}(volumen, balanceo de clases, dimensionalidad) hacen que los modelos basados en árboles de decisión superen a las redes neuronales en este dominio.
    \item \textbf{Evaluar el impacto del horizonte temporal} en la degradación del rendimiento predictivo de cada tipo de modelo.
    \item \textbf{Visualizar un datamart analítico} como parte de la infraestructura para integrar predicciones con datos históricos y facilitar la toma de decisiones    
\end{enumerate}

\section{Contribuciones de la Tesis}
Las principales contribuciones de este trabajo son:

\begin{enumerate}
    \item \textbf{Comparación rigurosa de modelo}: Se evaluaran tres enfoques de machine learning para predicción de riesgo crediticio multi-horizonte, demostrando que los modelos basados en árboles de decisión superan a las redes neuronales en datos tabulares con secuencias cortas.
    \item \textbf{Pipeline MLOps completo}: Se diseñó e implementó un sistema end-to-end que cubre desde la ingesta de datos hasta la visualización, con orquestación automatizada, gestión de experimentos y observabilidad.
    \item \textbf{Datamart analítico integrado}: Se creó un esquema estrella que permite cruzar datos históricos con predicciones del modelo, facilitando el análisis exploratorio y la toma de decisiones.
    \item \textbf{Dashboards interactivos}: Se implementaron visualizaciones en Apache Superset que permiten a los usuarios de negocio explorar datos y predicciones de forma autónoma.
\end{enumerate}

\section{Estructura del Documento}
El presente trabajo de titulación consta de los siguientes capítulos:

\begin{itemize}
    \item {Capitulo 1 (Introducción)}: Presenta el contexto, motivación, planteamiento del problema y objetivos de la investigación.
    \item {Capítulo 2 (Marco Teórico)}: Describe los fundamentos teóricos de inteligencia de negocio, aprendizaje automático y las tecnologías utilizadas.
    \item {Capítulo 3 (Estado del Arte)}: Revisa la literatura existente sobre predicción de riesgo crediticio, datamarts analíticos e integración de IA con plataformas BI.
    \item {Capítulo 4 (Metodología)}: Detalla el diseño del sistema, los datos utilizados, los modelos implementados y la configuración experimental.
    \item {Capítulo 5 (Resultados)}: Presenta los resultados experimentales, la comparación de modelos y el análisis de los hallazgos principales.
    \item {Capitulo 6 (Conclusiones) }: Expone el cumplimiento de los objetivos, las limitaciones del estudio y las recomendaciones para trabajos futuros.
    \item{Apéndice}: Listado de repositorios técnico del desarrollo del presente trabajo de titulación. 
\end{itemize}
%[Ov}Pasar a revision 